\section{Methodology}
\label{sec:methodology}

This section describes the evaluation protocol before the models. That
ordering is deliberate: the protocol is the contribution, and the models
are instruments for measuring what the protocol choices are worth. The
evaluation pipeline - the chronological split, the lagged/oracle
feature-regime fork, the daylight filter, and the choice of reference
forecast - is diagrammed in full in the repository (available at an
anonymised repository, link withheld for review).

\subsection{Forecast Task and Alignment}
\label{sec:methodology-task-alignment}

The task is deterministic point forecasting of hourly mean AC power at
horizons of 1, 3 and 6 hours ahead. Each horizon is treated as a
separate problem and reported separately; horizon dominates every other
factor in aggregate analyses of this literature \cite{nguyenmusgens2022meta}, so pooling across horizons would obscure the effects we set out to
measure.

The feature matrix is indexed by target time t. A row at target time t
for horizon h is the input to a forecast issued at time t-h. Every
observed quantity in that row must therefore originate at t-h or
earlier. Quantities describing t itself are admissible only when they
are deterministic and computable in advance: solar geometry, clear-sky
irradiance and power, and calendar encodings. This convention is
enforced for the tabular feature matrix by an assertion that
independently reconstructs each lagged column from the source series and
verifies the applied shift, together with a correlation probe that
flags any feature correlating implausibly with the target. For the
sequence inputs to the recurrent models, the assertion is additionally
validated by deliberate corruption: it is confirmed to fire both when
the window is shifted by one step and when the target-time value is
substituted into the window's final timestep, which is the specific
failure the convention exists to prevent.

\subsection{Evaluation Protocol}
\label{sec:eval-protocol}

We state the protocol as a set of explicit rules, in the spirit of the
model information sheets proposed by Kapoor and Narayanan
\cite{kapoor2023leakage}.

1. Splits are chronological and never shuffled. Training is 2011-2013,
   validation 2014, test 2015, each a whole calendar year.

2. Evaluation is restricted to daylight hours, defined as solar
   elevation above 10 degrees (80 degree zenith) - stricter than the 85
   degree zenith Yang et al. describe as typical practice
   \cite{yang2020verification}, because at low elevation the clear-sky
   index becomes unstable as measurement uncertainty and clear-sky
   model error grow against a small denominator. Night hours contain
   near-zero power that every model predicts correctly, and their
   inclusion is treated as a protocol variable rather than a default
   (Section \ref{sec:rq1}).

3. All scalers and feature statistics are fitted on training data only.
   The scaler implementation permits a single fit per instance and
   raises on a second call, so refitting on validation or test is not
   reachable by accident.

4. The headline metric is the root-mean-square-error skill score against
   the optimal convex combination of climatology and persistence,
   following the recommendation of Yang et al. \cite{yang2020verification}. The skill score
   against plain persistence is reported alongside, because the
   difference between the two is a result rather than a robustness
   check.

5. The lagged and oracle feature regimes are never mixed within a
   single model. Oracle features carry an explicit prefix, so that they
   cannot be mistaken for admissible predictors in a feature-importance
   analysis, and every oracle result is labelled a perfect-forecast
   upper bound.

6. The residual correction stage is fitted on out-of-fold residuals
   generated within the training period, described in Section IV-F.

7. Every experiment writes a machine-readable record containing its
   configuration, all metrics, timings, sample counts, random seed, the
   git commit hash, and a flag indicating whether the working tree was
   clean at execution time.

No hyperparameter tuning was performed, on any split, for any model;
the cost of this decision is noted in Section
\ref{sec:threats-to-validity}. Each rule above addresses a specific
leakage type in the taxonomy of Kapoor and Narayanan
\cite{kapoor2023leakage}, mapped in full in the repository.

\subsection{Feature Regimes}
\label{sec:methodology-feature-regimes}

Two feature regimes are constructed, and the distinction between them
is the second protocol axis we measure. The lagged regime contains 37
features in three groups: nine deterministic at target time (clear-sky
power and irradiance, solar geometry at the hour midpoint, cyclic
hour-of-day and day-of-year encodings), nineteen observations at or
before the issue time (measured power, clear-sky indices at the issue
time, two preceding hours, and a fixed 24-hour lag, five weather
channels, and staleness indicators), and nine rolling statistics whose
windows terminate at the issue time. The oracle regime adds five
measured weather channels at the target time, giving 42 features; it
is not a forecasting method but bounds what perfect weather knowledge
would be worth, the quantity against which architectural gains should
be judged (full feature list in the repository).

Two construction details matter. The daily lag is a fixed 24 hours at
every horizon, not a horizon-relative offset, so it stays aligned to
hour-of-day. And the clear-sky index is undefined at night, so at a
six-hour horizon - where a midday target is issued near dawn -
constructing lag features naively discards most midday targets:
daylight target retention was 75.6, 56.5 and 28.0 percent at horizons
of 1, 3 and 6 hours. We forward-fill the clear-sky indices with a
24-hour limit before shifting, add explicit staleness features, and
compute rolling statistics over the last N valid observations rather
than N wall-clock hours; retention is then 99.0 percent at all three
horizons.

\subsection{Reference Forecasts}
\label{sec:methodology-reference-forecasts}

Three reference forecasts are used.

Smart persistence forecasts
$\hat{P}(t) = k_p(t-h) \times P_{cs}(t)$, persisting the
clear-sky index from the issue time. Where no valid index exists at the
issue time it is forward-filled with a 24-hour limit, matching the
feature layer.

Climatology forecasts
$\hat{P}(t) = \bar{k}_p(\text{month}, \text{hour}) \times P_{cs}(t)$,
where the
mean clear-sky index per calendar month and hour of day is estimated on
training data only. It uses no recent observation, so the horizon does
not affect it.

The convex reference is w x persistence + (1-w) x climatology, with w
selected by grid search in steps of 0.01 to minimise root-mean-square
error on the validation split, fitted independently for each array,
horizon, seed and regime. Following Yang et al. \cite{yang2020verification}, this is the
reference for the headline skill score.

The third reference is necessary because the first degrades structurally
with horizon. At a six-hour horizon a midday target is issued near dawn,
so the persisted clear-sky index is stale: the proportion of daylight
forecasts relying on a forward-filled rather than directly observed
index rises from 3.3 percent at one hour to 22.6 percent at three hours
and 51 percent at six. The resulting bias reaches -2.15 kW at midday on
a 5 kW array. A skill score measured against a reference that fails in
this way rewards a model for handling staleness rather than for
forecasting, and the models can do so - they receive explicit staleness
features - while the reference cannot.

\subsection{Metrics}
\label{sec:methodology-metrics}

We report mean bias error, mean absolute error, root-mean-square error,
root-mean-square error normalised by nameplate capacity, and the skill
score $s = 1 - \mathrm{RMSE}_{forecast} / \mathrm{RMSE}_{reference}$.
Normalisation is by nameplate capacity rather than mean observed power,
so values are comparable across arrays of different sizes; normalised
error is not comparable across studies with different data regardless
of the normaliser \cite{yang2020verification}, and we make no
cross-study accuracy comparisons on that basis. All metrics for a
given cell are computed on the intersection of timestamps for which
every compared forecast produced a prediction, with documented outage
hours excluded first; that sample count is recorded with every result.

\subsection{Models}
\label{sec:methodology-models}

Five models are evaluated, chosen to decompose a hybrid architecture
rather than to survey the field. Gradient boosting over the tabular
feature matrix is the non-recurrent reference point; trees are
invariant to monotone feature transforms, so nothing is fit on training
data and nothing can leak. The recurrent model applies a single LSTM
layer over a 24-step sequence of nine per-hour channels ending at the
issue time, concatenates the final hidden state with the deterministic
target-time features, and maps the result through a two-layer head -
the deterministic features are necessary, since without them the
network has no representation of what time of day it is forecasting
for. The con\-vo\-lu\-tion\-al-recurrent model prepends a
one-dimensional convolution over the same sequence with same-length
padding, which does not leak because the window lies at or before the
issue time by construction.

The two residual-corrected models wrap either recurrent base with a
gradient-boosted stage fitted to the base model's residuals, generated
out of fold within the training period by an expanding window (fit on
year one, predict year two; fit on years one-two, predict year three;
residuals pooled), then the base model is refitted once on the full
training period, with validation used only for early stopping - each
fold's base model sees less data than the deployed one, so the pooled
residuals run slightly pessimistic, the correct direction to err.
Fitting the residual stage on validation residuals instead is reachable
behind an explicit flag that warns and records its use; it exists so we
can measure what it produces, not because we endorse it.

\subsection{Significance Testing}
\label{sec:methodology-significance-testing}

Model comparisons use the Diebold-Mariano test on paired squared-error
series \cite{diebold1995comparing}; because h-step-ahead errors are
autocorrelated by construction, the long-run variance uses a Bartlett
kernel truncated at lag h-1, with the small-sample correction of
Harvey, Leybourne and Newbold \cite{harvey1997testing} and
Holm-Bonferroni-adjusted p-values across the twenty-one pairwise
comparisons per array-horizon cell.

Seed variance across five seeds is reported separately and is not a
significance test: it measures whether a result recurs on retraining,
not whether one forecast beats another on this evaluation sample.
